\section{Conclusion}\label{sec:conclusion}

This paper studies a specific consequence of Bayesian parameter uncertainty: how much integrating a posterior for historical volatility changes an Average Price Option point value relative to plugging in the posterior mean. That mean-price correction is distinct from the width of the posterior distribution of theoretical prices and from the choice of volatility information supplied to the pricing model.

The synthetic evidence shows that the integration correction is not generically negligible. Across 175,000 simulated histories and 168 WTI-like contract states, the largest observed absolute PI--PM gap is 0.1384. The largest corrections occur with short historical samples, high volatility, long remaining horizons, extreme moneyness, and little realized fixing. In those states, the second-order approximation
\[
\frac12 C^{Q\prime\prime}(\bar\sigma)
\operatorname{Var}(\sigma\mid\mathcal D)
\]
tracks the realized integration gap closely. Posterior integration can therefore become appreciable when posterior dispersion and volatility curvature are jointly large.

The observed WTI application lies in a different region of that state space. The PI--PM correction is small relative to the width of the posterior distribution of model prices. Close PI and PM values therefore do not imply that parameter-induced price uncertainty is negligible; they imply that nonlinear averaging shifts the posterior mean price very little in these states. Independent pseudo-random Monte Carlo, randomized Sobol, and Curran calculations resolve this small correction well below its observed magnitude.

The larger empirical difference is between volatility inputs. Across 2,381 strict forward-in-time holdouts, prior-date APO inputs reduce pricing error sharply relative to the long-window historical-volatility benchmark. The evaluated rolling-window and exponentially weighted historical estimators do not reproduce that performance, and a horizon-matched Gaussian GARCH(1,1) forecast remains close to the long-window historical result. Reconstructing physical returns contract by contract from first-nearby CL settlements also leaves the comparison essentially unchanged. These checks narrow the claim to the tested historical specifications rather than all possible historical volatility models.

The persistence benchmark further narrows the interpretation of the APO result. On the 2,366 exact common rows, carrying the same contract's prior IV forward produces lower pooled error than the fitted previous-day smile. This descriptive ranking is not used to claim general superiority of carry; it establishes only that flexible cross-sectional smile fitting is unnecessary for the large historical-versus-option-informed improvement.

The independent vanilla-WTI experiment provides the cross-instrument check. A strictly prior LO surface materially improves on the historical benchmark on the common-support sample, and replacing the primary fixing-count RMS transfer with an arithmetic-average variance-matched scalar changes the result only slightly. The option-informed pricing advantage therefore survives both a different WTI option family and the alternative scalar transfer.

These results do not identify a pure structural difference between physical and risk-neutral volatility. Option-implied states can absorb risk premia, stochastic volatility, jumps, omitted futures-curve dynamics, liquidity, and market-marking effects. What the evidence establishes more narrowly is that, under the maintained one-factor pricing map, prior option-implied volatility states predict the observed APO settlements much better than the tested historical-volatility inputs, while posterior integration of the historical-volatility distribution changes the point-price mean only slightly.

Several limitations remain. The pricing dynamics are deliberately one-factor and constant-volatility. Discounting uses a Treasury par-yield approximation, and positive-volume observations remain sparse in the longest-dated APO maturities. The seven-expiry forward experiment contains 2,381 contract-date observations but only 15 distinct target dates, while the external LO common-support experiment contains 10 target dates. Valuation-date cluster resampling preserves contemporaneous dependence among contracts observed on the same date, but it does not create additional market histories or automatically remove serial dependence across adjacent dates. The external LO inversion is based on a 160-step American CRR tree; a stratified 80/160/320 refinement audit finds a median absolute IV change of 0.000284 from 160 to 320 steps, supporting practical numerical stability without implying exact convergence.

The paper's contribution is diagnostic rather than a new option-pricing theory. It separates the posterior-integration mean correction, posterior model-price dispersion, numerical error, historical-volatility recency, dynamic historical forecasting, option-state persistence, and volatility-input specification. The resulting benchmark is sharper than the original historical-versus-implied comparison: posterior integration becomes important when dispersion meets pricing curvature, but in the observed WTI APO sample the dominant empirical information is contained in prior option-implied volatility states rather than in either static or dynamically forecast historical volatility. Richer futures-curve or stochastic-volatility models can build on that benchmark without conflating model enrichment with the narrower effect of Bayesian posterior integration.
